\section{Determinism and Free Will in the History of Philosophy}

\subsection{The Myth}

\begin{itemize}
    \item  \textit{ก่อนหน้าที่มนุษย์จะพบคำตอบของปัญหาทั้งหลายจากศาสนาต่างๆ คำอธิบายทางศาสนาเหล่านี้ถูกถ่ายทอดต่อกันมาจากรุ่นหนึ่งสู่อีกรุ่นหนึ่งในรูปแบบของตำนานทวยเทพ (Myth) }
    \begin{itemize}
        \item จาก โลกของโซฟี่ หน้าที่ 22 
    \end{itemize}
\end{itemize}
     มนุษย์ในยุคนี้ใช้ Myth เป็นคำอธิบายของปัญหาต่างๆที่มนุษย์หาคำตอบไม่ได้, เป็นเครื่องมือในการส่งต่อความรู้,  และปกครองคนในกลุ่ม \\
\begin{itemize}
    \item \textit{เห็นได้จาก Edda เรื่อง “บทเพลงของทริม” ตอนที่กีริดื่นขึ้นมาแล้วพบว่าค้อนประจำตัวหายไป ธอร์โกรธจนมือสั่นและหวนคิดคราตั้ง เขาขวนก้องชื่อโลกิไปยืมปีกของเฟรย์จา เพื่อให้โลกิใส่บินไปยังโจตุนเฮม ดินแดนแห่งยักษ์ สืบดูว่าพวกมันขโมยค้อนของธอร์ไปหรือเปล่า }
    \item \textit{ที่โจตุนเฮม โลกิได้พบกับทริม ราชาแห่งยักษ์ ซึ่งมั่นอกมั่นใจมากจนอดคุยโม้ไม่ได้ว่ามันเอาค้อนไปซ่อนไว้ใต้พื้นโลกถึง 21 ไมล์ แถมยังบอกด้วยว่าธอร์ไม่มีทางจะได้ค้อนคืน นอกจากจะยอมให้มันแต่งงานกับเฟรย์จา} 
    \item \textit{ในทันใด เทพก็พบว่าพวกเขากำลังเผชิญกับวิกฤตการณ์ตัวประกัน ยักษ์ได้ลักเอาอาวุธที่ทรงอานุภาพที่สุดของพวกเทพไป ซึ่งเป็นสถานการณ์ที่ยอมไม่ได้เป็นอย่างยิ่ง เพราะตราบเท่าที่ยักษ์มีค้อนของธอร์อยู่ในมือ ก็เท่ากับพวกมันสามารถควบคุมโลกของเทพและมนุษย์ด้วย ถ้าต้องการค้อนไปก็ต้องแลกกับเฟรย์จา แต่นี่ก็เป็นข้อเสนอที่รับไม่ได้อีกเช่นกัน เพราะถ้านำเทพแห่งความอุดมสมบูรณ์ผู้ถูกเหล่าทวยเทพพิธีไปแลกกับยักษ์ หญ้าก็จะหายไปจากท้องทุ่ง และชีวิตของเทพและมนุษย์ก็จะถึงจุดจบด้วย}
    \item \textit{ตามตำนานโลกิกลับมายังแอสการ์ด และบอกให้เฟรย์จาเตรียมชุดเจ้าสาวเพราะเธอจะต้องวิวาห์กับราชายักษ์ เฟรย์จาโกรธเป็นฟืนเป็นไฟ และกล่าวว่าหากเธอยอมทำตาม พวกมนุษย์ก็จะว่าได้ว่าเธออยากมีสามีเต็มทีถึงยอมแต่งงานกับยักษ์}
    \item \textit{ทันใด เทพธอร์โมโหสุดขีดก็รีบไปหาแนะให้ธอร์แต่งตัวเป็นเจ้าสาวเสียเอง เพราะแค่เล่าผม แล้วเอาผ้าก้อนหินยัดใส่นอกเสื้อ ธอร์ก็จะดูเหมือนผู้หญิงแล้ว แม้ธอร์จะไม่เห็นด้วยกับแผนนี้เท่าไร ซึ่งก็พอเข้าใจได้ แต่ท้ายที่สุดก็ต้องยอม เพราะเป็นหนทางเดียวที่จะได้ค้อนคืนมา}
    \begin{itemize}
        \item จาก โลกของโซฟี่ หน้าที่ 25 
    \end{itemize}
\end{itemize}
สิ่งที่ Myth นี้กำลังอธิบาย 
\begin{itemize}
\item \textit{เมื่อเกิดความแล้ง มนุษย์พยายามหาคำอธิบายว่าเป็นไปได้มั้ยที่ ยักษ์ที่เป็นตัวแทนความชั่วร้าย ไปขโมยเครื่องมือที่ทำให้เกิดฝน (ค้อนของธอร์) และการที่ฝนกลับตกแสดงว่าค้อนของธอร์กลับมาแล้ว}\\
\begin{itemize}
        \item จาก โลกของโซฟี่ หน้าที่ 26
    \end{itemize}
 \end{itemize}
และผมว่าอีกเรื่องที่ Myth นี้แฝงไว้\\
ในสถานการณ์ที่ เราไม่มีค้อน(ทหาร)อำนาจต่อรอง \\
และต้องแลกเทพแห่งความอุดมสมบูรณ์ (น่าจะเป็นตัวแทนของประชาชนหรือเมือง) ไป\\
ผมว่าสิ่งธอร์ทำน่าจะเป็นเรื่องน่าอายในสมัยนั้น แปลว่า\\
ถ้าจะต้องแลกเมืองหรือ ประชาชน จะต้องทำเรื่องน่าอายขนาดไหนก็ต้องทำ\\
ขนาดธอร์ยังทำเลย
Myth ของคนในยุคนี้จึงไม่ใช่แค่เหล่าเพื่อความบรรเทิง แต่เป็นวิธีคิดและภูมิปัญญา
และ Myth ในยุคนี้มี Determinism แฝง เช่น
เทพพยากรณ์แก่งเดลฟี (โลกของโซฟี่ หน้าที่ 54) \\
ในความเห็นของผมในสมัยนั้น ชีวิตเราจะเป็นจะตายขึ้นอยู่กับธรรมชาติที่เทพสั่งในเป็น
เลยอาจจะเป็นสาเหตุที่ Determinism อยู่เต็ม Myth ในยุคนั้น



\subsection{The Natural Philosophers}
ปรัชญาสมัยใหม่ เช่น Descartes และ Kant

\subsection{The Athenian Philosophers}
ปรัชญาร่วมสมัย เช่น Existentialism และ Analytic Philosophy

\subsection{Hellenism}
ปรัชญาร่วมสมัย เช่น Existentialism และ Analytic Philosophy

\subsection{The Middle Ages}
ปรัชญาร่วมสมัย เช่น Existentialism และ Analytic Philosophy

\subsection{The Renaissance}
ปรัชญาร่วมสมัย เช่น Existentialism และ Analytic Philosophy

\subsection{The Baroque}
ปรัชญาร่วมสมัย เช่น Existentialism และ Analytic Philosophy

\subsection{The Early Modern Philosophy}
sdsadas

\subsection{The Enlightenment}
ปรัชญาร่วมสมัย เช่น Existentialism และ Analytic Philosophy

\subsection{The Modern}
ปรัชญาร่วมสมัย เช่น Existentialism และ Analytic Philosophy

\subsection{The Post Modern}
ปรัชญาร่วมสมัย เช่น Existentialism และ Analytic Philosophy